\section{Homotopy fixed points and Tate spectra}
\label{app:spectral-foundations}

In this appendix we compute the homotopy groups
of \eqref{eq:tate-compression}.
Let $R$ be a nonzero $p$-torsionfree perfectoid ring, with
$A,I_j,J_0^r,B,\Phi_B$ as in Section~\ref{sec:uniform}.
Every $I_j$ is a Cartier ideal.

\subsection{Homotopy groups of homotopy fixed-point and Tate spectra}

We first treat degrees that lift to virtual complex representations of
the circle. We use the same lift at each finite subgroup to pass from
the circle calculation to homotopy fixed points and Tate spectra for $C_{p^r}$.

\begin{definition}\label{def:app-local-spectra}
Let $U=\rho_p(\mathbf a)$ with $\mathbf a\in D_n^{(0)}$. The map
$\Sigma_n^{a_n}$ gives the identification
\[
E_0(U)\xrightarrow{\sim}B^{a_n}
        \prod_{j=1}^{n-1}(J_0^j)^{a_j-a_{j+1}}.
\]

Let $\mathbb C(k)$ be the complex $\mathbb T$-representation
of weight $k$. Put
\[
\widetilde U
=a_n\mathbb C(0)
+\sum_{j=0}^{n-1}(a_j-a_{j+1})\mathbb C(p^j).
\]
Its underlying real representation restricts to $U$.
For every $r$, we use the restriction of this same
$\widetilde U$ to $C_{p^r}$ to compute the homotopy groups.
Thus $E_0(U)$ is defined using $\widetilde U$, with $n$ fixed.

\end{definition}

By \cite[Section~2, p.~12]{SulymaNew},
$\THH(R\mid\mathbb S_{W(R^\flat)};\Z_p)\simeq X_R$
in $\Sp^{B\mathbb T}$.
Apply \cite[Theorem~1.9 and Remarks~3.1,~3.5]{SulymaNew}
with sequence $(a_0,\ldots,a_n;a_n)$.
After pullback along $\varphi$, using
$\varphi^*\operatorname{BK}(A,I_0)\cong B$, we obtain
\[
\begin{aligned}
\pi_0\bigl((S^{-\widetilde U}\wedge X_R)^{h\mathbb T}\bigr)
&\cong I_0^{a_0^+}E_0(U),\\
\pi_0\bigl((S^{-\widetilde U}\wedge X_R)^{t\mathbb T}\bigr)
&\cong E_0(U).
\end{aligned}
\]
These spectra have no odd homotopy groups, and the canonical map
is induced by $I_0^{a_0^+}E_0(U)\hookrightarrow E_0(U)$.

On restriction to $C_{p^r}$, the canonical map is still induced by this
inclusion, now on the quotients described below.

\begin{proposition}\label{prop:app-finite-comparison}
Let $U=\rho_p(\mathbf a)$, $\mathbf a\in D_n^{(0)}$,
with the $\mathbb T$-lift of Definition~\ref{def:app-local-spectra}.
For $r\geq0$ and $k\in\mathbb Z$,
\begin{align}
 \pi_0\bigl((S^{-U}\wedge X_R)^{hC_{p^r}}\bigr)&\cong
 \frac{I_0^{a_0^+}E_0(U)}{J_0^r I_0^{(a_0+1)^+}E_0(U)},
 &\pi_{2k+1}\bigl((S^{-U}\wedge X_R)^{hC_{p^r}}\bigr)&=0,
 \label{eq:app-finite-borel}\\
 \pi_0\bigl((S^{-U}\wedge X_R)^{tC_{p^r}}\bigr)&\cong \frac{E_0(U)}{J_0^rE_0(U)},
 &\pi_{2k+1}\bigl((S^{-U}\wedge X_R)^{tC_{p^r}}\bigr)&=0.
 \label{eq:app-finite-tate}
\end{align}
The canonical map from homotopy fixed points to the Tate spectrum is induced by
$I_0^{a_0^+}E_0(U)\hookrightarrow E_0(U)$.
\end{proposition}

\begin{proof}
Locally choose $I_0=(\xi)$ and $B=Ae$, with dual basis
$e^\vee\in B^\vee=\pi_{-2}X_R^{h\mathbb T}$.
By \cite[Lemma~3.1 and its proof]{Riggenbach}, for each $m\ge0$
we have the cofiber sequence
\[
\Sigma^{-2}M_m^{h\mathbb T}
\xrightarrow{\left(\prod_{j=1}^{r}\varphi^j(\xi)\right)e^\vee}
M_m^{h\mathbb T}
\xrightarrow{\operatorname{res}}
M_m^{hC_{p^r}},
\qquad
M_m=S^{m\mathbb C(1)-\widetilde U}\wedge X_R.
\]

For $K=\mathbb T,C_{p^r}$, the model
$EK=S(\mathbb C(1)^{\oplus\infty})|_K$ gives
\[
\operatorname*{colim}_{m\ge0}M_m^{hK}
\simeq
\bigl(S^{-\widetilde U}\wedge X_R\bigr)^{tK},
\]
where the transition maps are induced by
$S^0\to S^{\mathbb C(1)}$.
Under this equivalence, the map from $m=0$ is
$\operatorname{can}$.
Taking colimits of the cofiber sequences above gives the
bottom row of the following commutative diagram;
the top row is the sequence at $m=0$:
\[
\begin{tikzcd}[column sep=large,row sep=huge]
\Sigma^{-2}(S^{-\widetilde U}\wedge X_R)^{h\mathbb T}
  \arrow[r,"{\left(\prod_{j=1}^r\varphi^j(\xi)\right)e^\vee}"]
  \arrow[d,"\Sigma^{-2}\operatorname{can}"']
&
(S^{-\widetilde U}\wedge X_R)^{h\mathbb T}
  \arrow[r,"\res"]
  \arrow[d,"\operatorname{can}"]
&
(S^{-\widetilde U}\wedge X_R)^{hC_{p^r}}
  \arrow[d,"\operatorname{can}"]
\\
\Sigma^{-2}(S^{-\widetilde U}\wedge X_R)^{t\mathbb T}
  \arrow[r,"{\left(\prod_{j=1}^r\varphi^j(\xi)\right)
    \operatorname{can}(e^\vee)}"']
&
(S^{-\widetilde U}\wedge X_R)^{t\mathbb T}
  \arrow[r]
&
(S^{-\widetilde U}\wedge X_R)^{tC_{p^r}}.
\end{tikzcd}
\]

For $k\in\mathbb Z$, the left square on $\pi_{2k}$ is
\[
\begin{tikzcd}[column sep=large,row sep=large]
I_0^{(a_0+k+1)^+}B^{k+1}\otimes_A E_0(U)
  \arrow[r,hook]
  \arrow[d,hook]
&
I_0^{(a_0+k)^+}B^k\otimes_A E_0(U)
  \arrow[d,hook]
\\
B^{k+1}\otimes_A E_0(U)
  \arrow[r,hook]
&
B^k\otimes_A E_0(U),
\end{tikzcd}
\]
with horizontal maps
\[
ae^{k+1}\otimes\ell
\longmapsto
\left(\prod_{j=1}^r\varphi^j(\xi)\right)
ae^k\otimes\ell.
\]
Since
\[
J_0^r=\left(\prod_{j=1}^r\varphi^j(\xi)\right),
\]
the long exact sequences give
\eqref{eq:app-finite-borel}--\eqref{eq:app-finite-tate}.
The right vertical inclusions induce
\[
\operatorname{can}:
\frac{I_0^{a_0^+}E_0(U)}
     {J_0^rI_0^{(a_0+1)^+}E_0(U)}
\longrightarrow
\frac{E_0(U)}{J_0^rE_0(U)},
\qquad
[x]\longmapsto[x].
\]
\end{proof}

\subsection{Maps from geometric fixed points to Tate spectra}

We next compute the map from geometric fixed points to Tate spectra.
For $r\ge1$ and $U\in RO(C_{p^r})$, let
\[
\theta_U:
\Phi^{C_p}(S^{-U}\wedge X_R)
\longrightarrow (S^{-U}\wedge X_R)^{tC_p}
\]
be the canonical map. Let
$q:C_{p^r}\to C_{p^r}/C_p$ and write
$U=q^*W+M$, with $W\in RO(C_{p^r}/C_p)$ and $M^{C_p}=0$.
For actual $M$, the sphere $S(M)$ is $C_p$-free, so the Euler map
$a_M$ is a Tate equivalence and $\Phi^{C_p}(a_M)=\mathrm{id}$.
For virtual $M$, it's implied by inverse Euler equivalences for negative summands.
By naturality in \cite[Lemma~II.4.5]{NS}, we have
\begin{equation}\label{eq:app-moving-arrow}
\theta_U=(a_M^{tC_p})^{-1}\circ(S^{-W}\wedge\theta_0).
\end{equation}
Take $r=n\ge1$ and
$U=\rho_p(\mathbf a)$ with $\mathbf a\in D_n^{(0)}$.
Under the cyclotomic equivalence, $\theta_0$ is cyclotomic
Frobenius. Put
\[
W=\rho_{C_{p^{n-1}}}(a_1,\ldots,a_n),
\qquad U=q^*W+(a_0-a_1)\lambda_0.
\]
Definitions~\ref{def:natural-degree}
and~\ref{def:app-local-spectra} give
\begin{align}
\varphi^*E_0(W)
  &=I_1^{-a_1}E_0(U),\notag\\
\varphi^*\bigl(I_0^{a_1^+}E_0(W)\bigr)
  &=I_1^{(-a_1)^+}E_0(U).
\label{eq:app-frob-numerator}
\end{align}

For the circle lift of Definition~\ref{def:app-local-spectra},
the Frobenius arrow in \cite[Theorem~1.9, Notation~3.14 and Example~3.15]{SulymaNew},
using the product of fractional ideals for virtual weights,
with the Breuil--Kisin Frobenius of Subsection~\ref{subsec:period-line},
has $A$-linearization
\[
I_1^{(-a_1)^+}E_0(U)\hookrightarrow E_0(U)
\]
under \eqref{eq:app-frob-numerator}.

By \cite[Lemma~II.4.1]{NS}, restriction to
$C_{p^{n-1}}\le\mathbb T/C_p$ gives
\[
(S^{-W}\wedge X_R)^{hC_{p^{n-1}}}
\xrightarrow{\ \theta_U^{hC_{p^{n-1}}}\ }
(S^{-U}\wedge X_R)^{tC_{p^n}}.
\]
Since $\varphi(J_0^{n-1})=I_1^{-1}J_0^n$,
\begin{equation}\label{eq:app-frob-denominator}
\varphi^*\bigl(J_0^{n-1}I_0^{(a_1+1)^+}E_0(W)\bigr)
=J_0^n I_1^{(-a_1-1)^+}E_0(U)
\subseteq J_0^nE_0(U).
\end{equation}
Naturality and Proposition~\ref{prop:app-finite-comparison}
identify the induced finite quotient map with $[x]\mapsto[x]$.
Applying this calculation to $U^{C_{p^i}}$ at $C_{p^{n-i}}$
and pulling back along $\varphi^i$ gives $\ell_i$ and $\rho_i$
as the maps $[x]\mapsto[x]$ of Definition~\ref{def:natural-degree}.
For $i=n-1$, Proposition~\ref{prop:app-finite-comparison}
with $r=0$ gives
$\varphi^{n*}\pi_{2a_n}X_R=A_n(U)$.

\subsection{The sign representation at \texorpdfstring{$p=2$}{p=2}}

What remains is the selected natural degrees with an odd
coefficient of the sign representation. By the sign Euler cofiber sequence, it suffices to compute the local terms' kernels of restriction in complex degrees.
Assume $p=2$ and $n\geq1$. Let $U$ be a
selected natural degree with $\epsilon(\mathbf a)=1$, and put
$W=U-\sigma_{2^n}$. Then $W$ is complex, with degree coordinates
$(a_0-\tfrac12,\ldots,a_{n-1}-\tfrac12,a_n)$.
For $0\le i<n$, put $r=n-i$. Dualizing
\[
(C_{2^r}/C_{2^{r-1}})_+\longrightarrow S^0
\longrightarrow S^{\sigma_{2^r}}
\]
and applying Shapiro identifies the local terms for $U^{C_{2^i}}$
with the fibers of restriction from $C_{2^r}$ to $C_{2^{r-1}}$
for $W^{C_{2^i}}$, for both homotopy fixed points and Tate spectra.
Use the same circle lift at both subgroups.
By Proposition~\ref{prop:app-finite-comparison}, these restrictions
are surjective quotient maps in even degrees, so the fibers have
no odd homotopy groups.
In degree zero, we have
\begin{align}
 \pi_0\bigl((S^{-U^{C_{2^i}}}\wedge X_R)^{hC_{2^r}}\bigr)
   &\cong\frac{J_0^{r-1}I_0^{(a_i+1/2)^+}E_0(W^{C_{2^i}})}
           {J_0^rI_0^{(a_i+1/2)^+}E_0(W^{C_{2^i}})},\notag\\
 \pi_0\bigl((S^{-U^{C_{2^i}}}\wedge X_R)^{tC_{2^r}}\bigr)
   &\cong\frac{J_0^{r-1}E_0(W^{C_{2^i}})}{J_0^rE_0(W^{C_{2^i}})}.
 \label{eq:app-sign-local-kernels}
\end{align}

After pulling back along $\varphi^i$,
\[
J_i^{n-1}E_i(W)=\mathcal S_i(U),
\]
so \eqref{eq:app-sign-local-kernels} gives \eqref{eq:uniform-sign}.

Restrict $\ell_i,\rho_i$ for $W$ to these kernels.
If $r\ge2$, then $C_2\subseteq\ker(\sigma_{2^r})$, so
$(\sigma_{2^r})^{C_2}$ is the sign representation of $C_{2^r}/C_2$.
If $r=1$, then $(\sigma_{2^r})^{C_2}=0$ and $\Phi^{C_2}(a_{\sigma_2})$
is the identity. The map $\rho_{n-1}$ follows by naturality
of the sign Euler cofiber sequence.

